\section{Conclusions}
\label{sec:conclusions}

The Fermi paradox is usually modelled by varying how many civilizations arise. We
have argued that it should also be modelled by varying how long they remain
externally observable, and that observability is best treated as a state variable
with its own equation of motion rather than as a coefficient.

Three results follow.

First, once existence, signal production and search coverage are separated in
Eq.~\eqref{eq:decomp}, non-observation constrains their product and not $\Ntrue$
alone. This is a logical point, but it is routinely elided, and it determines
what a model of the paradox should vary. \Established

Second, an observability equation must be constructed so that it cannot leave the
physical range. Writing suppression as a fractional rate acting on emitted signal,
above an explicit thermodynamic floor, achieves this by construction
(Proposition~\ref{prop:obs}); the analogous treatment of erosion confines
regulatory capacity to the unit interval (Proposition~\ref{prop:reg}). These are
not cosmetic. The uncorrected formulation produced, and confidently labelled, a
``quiet'' civilization whose observability was $-1431$.

Third --- and this is the modelling result we would defend most strongly ---
peak-then-decline observability need not be assumed. Because the quasi-steady
observability is a ratio of production to suppression, the decline emerges
whenever suppression scales faster in capability than production does
(Eq.~\eqref{eq:decline}). A low-observability account of the Great Silence built
this way does not assume its conclusion, and can be attacked on the plausibility
of a scaling separation rather than on the choice of a curve. \Plausible

We have also reported a negative result. The Lambert $W$ threshold that named and
motivated this work is not a property of the integrated system: the natural
balance cancels $\Icap$ and yields a logarithm, and recovering Lambert structure
requires an auxiliary balance we must declare rather than derive
(\S\ref{sec:lambert}). We keep the boundary as a diagnostic and depend on it
nowhere.

Finally, the framework is verified rather than merely implemented. Closed-form
solutions, an independent high-order integrator, property-based testing over 300
randomised parameter sets, and two independently written implementations agreeing
to solver tolerance together establish that the equations described here are the
equations solved. They establish nothing about whether those equations describe a
real civilization, and we have been careful to keep verification and validation
apart.

The framework does not solve the Fermi paradox and is not intended to. Its
purpose is to make a class of arguments about cosmic silence precise enough to be
wrong: to say which assumptions produce which regimes, which conclusions survive
a change of functional form, and what observation would count against them. The
most useful next steps are a population-level treatment that can be set against
grabby-aliens constraints, a collapse term that produces trajectories rather than
flags, and a nondimensionalisation that reduces the parameter count to the few
groups that actually govern the behaviour.
